% Generated by tools/build_new_dists_anthology_sections.py; do not edit by hand.

\section{Multiplicative wreath representations}\label{proof:new-dist:multiplicative_wreath}

\resultbox{\textbf{Canonical testing artifacts.}\par\smallskip Exact backend: \path{new_dists/multiplicative_wreath_sigma_mgfs/verification.py}.\par Regression: \path{new_dists/multiplicative_wreath_sigma_mgfs/test_verification.py}.\par Marimo: \path{marimo_notebooks/multiplicative_wreath.py}.}

\subsection*{Scope and outcome}

Let $G_n=H\wr S_n$.  We prove the wreath class-power theorem and use it to
derive exact $\sigma$-MGFs for product actions on $\Omega^n$ and
tensor-induced modules $V^{\otimes n}$.  Their tensor moments are multiset
counts, which normally diverge polynomially; diagonal coset and twisted
regular actions give related polynomial or exponential rare-event laws.  We
also record the graded bosonic and fermionic Fock formulas, fixed-tail
multipartition limits, growing-multipartition obstructions, varying-base
criteria, and the contrasting compound-Poisson law for additive block
modules.  An executable marimo notebook constructs the relevant permutation
and signed-monomial matrix groups in dimensions up to $36$.  Direct orbit and
fixed-space computations agree with the proposed formulas in $13{,}183$
exact checks.

\subsection{Common target and wreath class powers}

Let $H$ be a finite group and $G_n=H^n\rtimes S_n$.  For a complex
$G_n$-module $W$ and a partition $\tau=(\tau_1,\ldots,\tau_s)$, put
\[
 a_\tau(G_n,W)=\dim\left(\bigotimes_{j=1}^s\Sym^{\tau_j}W\right)^{G_n},
 \qquad
 \MGF_{G_n,W}(\mathbf t)=\sum_\tau a_\tau m_\tau(\mathbf t).
\]
Reynolds averaging and
$\sum_{d\ge0}\Tr(\Sym^d A)t^d=\det(I-tA)^{-1}$ give
\begin{equation}\label{nd:multiplicative_wreath:eq:universal-molien}
 \MGF_{G_n,W}
 =\frac1{|G_n|}\sum_{g\in G_n}
 \exp\!\left(\sum_{r\ge1}\frac{p_r(\mathbf t)}r\chi_W(g^r)\right).
\end{equation}

Let $\mathcal C(H)$ be the conjugacy classes of $H$.  A wreath type is a
partition-valued function $\boldsymbol\mu:C\mapsto\mu(C)$ of total size $n$.
A part $\ell$ of $\mu(C)$ denotes an $\ell$-cycle of the top permutation whose
ordered label product lies in $C$.  For $c\in C$, set $z_C=|C_H(c)|$, and set
\[
 z_\lambda=\prod_{j\ge1}j^{m_j(\lambda)}m_j(\lambda)!,\qquad
 Z_{\boldsymbol\mu}=\prod_{C\in\mathcal C(H)}
 z_{\mu(C)}z_C^{\ell(\mu(C))}.
\]

\begin{theorem}[wreath class-power theorem]\label{nd:multiplicative_wreath:thm:class-power}
The proportion of elements of type $\boldsymbol\mu$ is
$Z_{\boldsymbol\mu}^{-1}$.  If an $\ell$-cycle has cycle product $c$ and
$d=\gcd(\ell,r)$, then the $r$th power has $d$ cycles of length $\ell/d$,
each with cycle product conjugate to $c^{r/d}$.  Denote the transformed type
by $\boldsymbol\mu^{[r]}$.  Then
\begin{equation}\label{nd:multiplicative_wreath:eq:universal-class}
 \boxed{\MGF_{G_n,W}=
 \sum_{\boldsymbol\mu\vdash_H n}\frac1{Z_{\boldsymbol\mu}}
 \exp\!\left[\sum_{r\ge1}\frac{p_r}{r}
 \chi_W(\boldsymbol\mu^{[r]})\right].}
\end{equation}
\end{theorem}

\begin{proof}
Fix the top cycle structure and one cycle of length $\ell$.  After choosing
$\ell-1$ labels freely, the last label has exactly $|C|$ choices that place
the cycle product in $C$.  Combining this with the ordinary permutation
cycle-index count gives the denominator $z_{\mu(C)}z_C^{\ell(\mu(C))}$;
multiplication over $C$ gives $Z_{\boldsymbol\mu}$.  Advancing by $r$ positions
on an $\ell$-cycle has $d=\gcd(\ell,r)$ orbits, each of length $\ell/d$.
Traversing one new orbit winds $r/d$ times around the old cycle, so its label
product is conjugate to $c^{r/d}$.  Substitution in
\eqref{nd:multiplicative_wreath:eq:universal-molien} proves \eqref{nd:multiplicative_wreath:eq:universal-class}.
\end{proof}

\begin{remark}[orientation]
Changing the starting point of a top cycle conjugates its label product.
Thus all formulas are independent of left-versus-right cycle notation once
the product is recorded only up to conjugacy.
\end{remark}

\subsection{Product-action permutation groups}

Suppose $H$ acts on a finite set $\Omega$ and $G_n$ acts on
$X_n=\Omega^n$.  Write $g=(h_1,\ldots,h_n;\pi)$.  For a cycle $Q$ of $\pi$,
let $a_Q$ be its ordered label product and put
$d_Q(r)=\gcd(|Q|,r)$ and $f_s(a)=|\Fix_\Omega(a^s)|$.

\begin{theorem}[product-action fixed powers and $\sigma$-MGF]
For every $r\ge1$,
\begin{equation}\label{nd:multiplicative_wreath:eq:product-fixed}
 \boxed{F_r(g):=|\Fix_{\Omega^n}(g^r)|
 =\prod_{Q\in\Cyc(\pi)}
 f_{r/d_Q(r)}(a_Q)^{d_Q(r)}.}
\end{equation}
If $c_j(g)$ is the number of $j$-cycles on $X_n$, then
\begin{equation}\label{nd:multiplicative_wreath:eq:mobius}
 c_j(g)=\frac1j\sum_{r\mid j}\mu(j/r)F_r(g),
\end{equation}
and therefore
\begin{equation}\label{nd:multiplicative_wreath:eq:product-smg}
 \boxed{\MGF^{\rm prod}_n
 =\sum_{\boldsymbol\mu\vdash_H n}\frac1{Z_{\boldsymbol\mu}}
 \prod_{i,j\ge1}(1-t_i^j)^{-c_j(\boldsymbol\mu)}.}
\end{equation}
\end{theorem}

\begin{proof}
The $r$th power splits $Q$ into $d_Q(r)$ coordinate cycles.  Going around one
new cycle applies an element conjugate to $a_Q^{r/d_Q(r)}$.  Each new cycle
may independently receive any point fixed by that element, proving
\eqref{nd:multiplicative_wreath:eq:product-fixed}.  Since
$F_r(g)=\sum_{j\mid r}j c_j(g)$, M\"obius inversion gives
\eqref{nd:multiplicative_wreath:eq:mobius}.  The represented permutation matrix has determinant
$\prod_j(1-t^j)^{c_j(g)}$; Reynolds averaging and
Theorem~\ref{nd:multiplicative_wreath:thm:class-power} give \eqref{nd:multiplicative_wreath:eq:product-smg}.
\end{proof}

Define $R_s(H,\Omega)=\#(H\backslash\Omega^s)$.

\begin{theorem}[exact product moments]\label{nd:multiplicative_wreath:thm:product-moment}
\begin{equation}\label{nd:multiplicative_wreath:eq:product-moment}
 \boxed{[m_{(1^s)}]\MGF^{\rm prod}_n
 =\binom{n+R_s(H,\Omega)-1}{R_s(H,\Omega)-1}.}
\end{equation}
\end{theorem}

\begin{proof}
An ordered $s$-tuple of words is an $s\times n$ array.  The base group $H^n$
replaces each column by its $H$-orbit in $\Omega^s$, leaving $R_s$ column
types.  The top group forgets their order.  The orbits are therefore
multisets of size $n$ from $R_s$ types, counted by \eqref{nd:multiplicative_wreath:eq:product-moment}.
The orbit count equals the invariant dimension in the permutation module.
\end{proof}

If the base action is nontrivial and transitive, then $R_2\ge2$, so the pair
coefficient grows as $n^{R_2-1}/(R_2-1)!$ and no finite coefficientwise limit
exists.  If $\delta_H$ is the derangement proportion of the base action, a
fixed word exists exactly when every top-cycle product is a nonderangement.
Since the cycle products are independent uniform elements of $H$ conditional
on $\pi$,
\begin{equation}\label{nd:multiplicative_wreath:eq:product-rare}
 \Pr(F_1>0)=\E(1-\delta_H)^{K(\pi)}
 =\frac{\Gamma(n+1-\delta_H)}
 {\Gamma(1-\delta_H)\Gamma(n+1)}
 \sim\frac{n^{-\delta_H}}{\Gamma(1-\delta_H)}.
\end{equation}
Thus $F_1\to0$ in probability while transitivity gives $\E F_1=1$ and higher
moments grow polynomially.

\subsection{Tensor-induced signed matrix groups}

Let $V$ be an $H$-module with character $\chi$, and let
$T_n(V)=V^{\otimes n}$ with the usual label action and permutation of tensor
factors.

\begin{theorem}[tensor cycle character and moments]\label{nd:multiplicative_wreath:thm:tensor}
For $g=(h_1,\ldots,h_n;\pi)$,
\begin{align}
 \chi_{T_n(V)}(g)&=\prod_{Q\in\Cyc(\pi)}\chi(a_Q),\label{nd:multiplicative_wreath:eq:tensor-char}\\
 \chi_{T_n(V)}(g^r)&=\prod_Q
 \chi(a_Q^{r/d_Q(r)})^{d_Q(r)}.\label{nd:multiplicative_wreath:eq:tensor-power}
\end{align}
Consequently,
\begin{equation}\label{nd:multiplicative_wreath:eq:tensor-smg}
 \boxed{\MGF_n^\otimes=
 \sum_{\boldsymbol\mu\vdash_Hn}\frac1{Z_{\boldsymbol\mu}}
 \exp\!\left[\sum_{r\ge1}\frac{p_r}{r}
 \prod_{C,\ell}\chi(c^{r/d_{\ell,r}})^{
 d_{\ell,r}m_\ell(\mu(C))}\right].}
\end{equation}
If $a_s(H,V)=\dim(V^{\otimes s})^H$, then
\begin{equation}\label{nd:multiplicative_wreath:eq:tensor-moment}
 \boxed{[m_{(1^s)}]\MGF_n^\otimes=
 \begin{cases}
 \binom{n+a_s-1}{a_s-1},&a_s>0,\\
 0,&a_s=0.
 \end{cases}}
\end{equation}
\end{theorem}

\begin{proof}
Contracting tensor indices around one top cycle produces the trace of its
label product, proving \eqref{nd:multiplicative_wreath:eq:tensor-char}; Theorem~\ref{nd:multiplicative_wreath:thm:class-power}
gives \eqref{nd:multiplicative_wreath:eq:tensor-power}.  Substitution in
\eqref{nd:multiplicative_wreath:eq:universal-class} proves \eqref{nd:multiplicative_wreath:eq:tensor-smg}.  For the moment,
rearrange tensor factors:
\[
 (V^{\otimes n})^{\otimes s}\cong(V^{\otimes s})^{\otimes n}.
\]
Taking $H^n$-invariants gives $((V^{\otimes s})^H)^{\otimes n}$; taking the
remaining $S_n$-invariants gives
$\Sym^n((V^{\otimes s})^H)$.  Its dimension is
\eqref{nd:multiplicative_wreath:eq:tensor-moment}.
\end{proof}

If $a_s\ge2$, the coefficient is asymptotic to
$n^{a_s-1}/(a_s-1)!$.  If $V$ is self-dual irreducible and $\dim V>1$, then
$a_2=1$ and
\[
 a_4=\dim\operatorname{End}_H(V\otimes V)\ge2,
\]
because $V\otimes V=\Sym^2V\oplus\Lambda^2V$ has two nonzero invariant
summands.  Hence the fourth tensor moment is at least $n+1$.  A
one-dimensional character $\theta$ of order $q$ is exceptional:
$a_s=1$ when $q\mid s$ and $a_s=0$ otherwise, so these obstruction
coefficients remain bounded.

\subsection{Bosonic and fermionic Fock modules}

For a finite $G$-module $W$, introduce the energy-graded modules
\[
 \mathcal F^+(W)=\bigoplus_{d\ge0}q^d\Sym^dW,
 \qquad
 \mathcal F^-(W)=\bigoplus_{d\ge0}q^d\Lambda^dW.
\]
The graded $\sigma$-MGF applies the outer symmetric-power construction to
these graded modules.

\begin{proposition}[graded Fock formulas]\label{nd:multiplicative_wreath:prop:fock}
\begin{align}
 \mathcal S^+_{G,W}(q;\mathbf t)
 &=\frac1{|G|}\sum_{g\in G}\exp\!\left[
 \sum_{r\ge1}\frac{p_r}{r}\det(I-q^rg^r\mid W)^{-1}\right],
 \label{nd:multiplicative_wreath:eq:fock-plus}\\
 \mathcal S^-_{G,W}(q;\mathbf t)
 &=\frac1{|G|}\sum_{g\in G}\exp\!\left[
 \sum_{r\ge1}\frac{p_r}{r}\det(I+q^rg^r\mid W)\right].
 \label{nd:multiplicative_wreath:eq:fock-minus}
\end{align}
Removing the vacuum replaces the determinant inverse in
\eqref{nd:multiplicative_wreath:eq:fock-plus} by that inverse minus $1$.
\end{proposition}

\begin{proof}
For each $r$, the graded characters are
\[
 \sum_{d\ge0}q^{dr}\Tr(g^r\mid\Sym^dW)
 =\det(I-q^rg^r\mid W)^{-1},
\]
and
\[
 \sum_{d\ge0}q^{dr}\Tr(g^r\mid\Lambda^dW)
 =\det(I+q^rg^r\mid W).
\]
Insert these identities into the plethystic form of Reynolds averaging.
\end{proof}

For the block module $W_n=V^{\oplus n}$, every total-energy-$E$ monomial has
support on at most $E$ blocks.  Therefore
\begin{equation}\label{nd:multiplicative_wreath:eq:fock-stability}
 [q^Em_\tau]\mathcal S^\pm_{H\wr S_n,W_n}
 \quad\text{is independent of $n$ for $n\ge E$.}
\end{equation}
The grading is essential: at $q=1$ the bosonic Fock space is
infinite-dimensional and ungraded coefficients are generally infinite.

\subsection{Irreducible multipartitions}

Irreducibles of $H\wr S_n$ are indexed by multipartitions
$\boldsymbol\lambda=(\lambda^\rho)_{\rho\in\Irr(H)}$.  The wreath Frobenius
characteristic sends their characters to
\begin{equation}\label{nd:multiplicative_wreath:eq:wreath-frob}
 \chi^{\boldsymbol\lambda}\longleftrightarrow
 \prod_{\rho\in\Irr(H)}s_{\lambda^\rho}(x^{(\rho)}).
\end{equation}
Theorem~\ref{nd:multiplicative_wreath:thm:class-power} immediately gives
\begin{equation}\label{nd:multiplicative_wreath:eq:irrep-smg}
 \boxed{\MGF_{\boldsymbol\lambda}
 =\sum_{\boldsymbol\mu\vdash_Hn}\frac1{Z_{\boldsymbol\mu}}
 \exp\!\left[\sum_{r\ge1}\frac{p_r}{r}
 \chi^{\boldsymbol\lambda}(\boldsymbol\mu^{[r]})\right].}
\end{equation}

For a fixed-tail family, only the trivial-color component acquires a long
first row and all other boxes form a fixed multipartition
$\boldsymbol\alpha$.  In the stable range the character is a polynomial
$Q_{\boldsymbol\alpha}(X_{\ell,C})$ in marked cycle counts.  Under the
uniform measure,
\begin{equation}\label{nd:multiplicative_wreath:eq:marked-poisson}
 X_{\ell,C}\Longrightarrow Z_{\ell,C},\qquad
 Z_{\ell,C}\sim\operatorname{Poisson}
 \left(\frac1{\ell|C_H(c)|}\right),
\end{equation}
independently.  Define
\begin{equation}\label{nd:multiplicative_wreath:eq:power-poisson}
 Z_{j,D}^{[r]}=
 \sum_{\substack{\ell,C:\ \ell/\gcd(\ell,r)=j\\
 c^{r/\gcd(\ell,r)}\in D}}
 \gcd(\ell,r)Z_{\ell,C}.
\end{equation}
Coefficientwise passage to the limit in \eqref{nd:multiplicative_wreath:eq:irrep-smg} yields the
marked-Poisson-chaos law
\begin{equation}\label{nd:multiplicative_wreath:eq:fixed-tail-limit}
 \boxed{\MGF_{\boldsymbol\alpha,\infty}
 =\E\exp\!\left[\sum_{r\ge1}\frac{p_r}{r}
 Q_{\boldsymbol\alpha}(\boldsymbol Z^{[r]})\right].}
\end{equation}

There is no corresponding universal raw limit for growing multipartitions.
If a self-dual $W_{\boldsymbol\lambda}$ satisfies
\[
 W_{\boldsymbol\lambda}\otimes W_{\boldsymbol\lambda}
 =\bigoplus_{\boldsymbol\nu}
 g^H_{\boldsymbol\lambda,\boldsymbol\lambda,\boldsymbol\nu}
 W_{\boldsymbol\nu},
\]
then
\begin{equation}\label{nd:multiplicative_wreath:eq:kronecker-obstruction}
 [m_{(1,1,1,1)}]\MGF
 =\dim\operatorname{End}_{H\wr S_n}(W_{\boldsymbol\lambda}^{\otimes2})
 =\sum_{\boldsymbol\nu}
 (g^H_{\boldsymbol\lambda,\boldsymbol\lambda,\boldsymbol\nu})^2.
\end{equation}
Growth of this sum obstructs convergence.

\subsection{Diagonal coset matrix groups}

Let $N_n=H^n$, $\Delta H=\{(h,\ldots,h):h\in H\}$, and let $N_n$ act on
$X_n=N_n/\Delta H$.  For a class $C$ and $c\in C$, write
$z_C=|C_H(c)|$.

\begin{theorem}[diagonal fixed powers and moments]\label{nd:multiplicative_wreath:thm:diagonal}
For $g=(g_1,\ldots,g_n)$,
\begin{equation}\label{nd:multiplicative_wreath:eq:diagonal-fixed}
 F_r(g)=\begin{cases}
 z_C^{n-1},&g_1^r,\ldots,g_n^r\in C,\\
 0,&\text{otherwise}.
 \end{cases}
\end{equation}
If $X(g)=F_1(g)$, then
\begin{equation}\label{nd:multiplicative_wreath:eq:diagonal-moment}
 \boxed{\E X(g)^s=\sum_{C\in\mathcal C(H)}z_C^{(s-1)n-s}.}
\end{equation}
In particular, $[m_{(1,1)}]=\sum_Cz_C^{n-2}\ge|H|^{n-2}$.
\end{theorem}

\begin{proof}
A coset $x\Delta H$ is fixed by $g^r$ exactly when the components of
$x^{-1}g^rx$ are equal.  If the component powers do not lie in one class,
there is no solution.  If they lie in $C$, choose the common conjugate and
then choose a conjugator in each coordinate; division by $|\Delta H|$ leaves
$z_C^{n-1}$ cosets.  The event $g_1,\ldots,g_n\in C$ has probability
$z_C^{-n}$, and the fixed-point value there is $z_C^{n-1}$.  Summing the
$s$th powers proves \eqref{nd:multiplicative_wreath:eq:diagonal-moment}.
\end{proof}

The identity class alone gives exponential second-moment growth.  Meanwhile
$\Pr(X>0)=\sum_Cz_C^{-n}\to0$ for nontrivial $H$, although $\E X=1$.

\subsection{Twisted-wreath regular matrix groups}

Let $N_n=H^n$, let $P_n\le\operatorname{Aut}(N_n)$, and let
$G_n=N_n\rtimes P_n$ act affinely on $N_n$ by
$x^{(a,p)}=a\,p(x)$.  Put
\[
 a_r=a\,p(a)\cdots p^{r-1}(a),
 \qquad L_{p^r}(x)=x\,p^r(x)^{-1}.
\]

\begin{theorem}[Lang-map criterion and moments]\label{nd:multiplicative_wreath:thm:twisted}
\begin{equation}\label{nd:multiplicative_wreath:eq:lang}
 F_r(a,p)=\begin{cases}
 |\Fix_{N_n}(p^r)|,&a_r\in\im L_{p^r},\\
 0,&\text{otherwise}.
 \end{cases}
\end{equation}
Conditioned on $p$, $F_1$ equals $|\Fix_{N_n}(p)|$ with probability its
reciprocal, and otherwise equals zero.  Hence
\begin{equation}\label{nd:multiplicative_wreath:eq:twisted-moment}
 \boxed{[m_{(1^s)}]=\frac1{|P_n|}\sum_{p\in P_n}
 |\Fix_{N_n}(p)|^{s-1}=\#(P_n\backslash N_n^{s-1}).}
\end{equation}
If $P_n=S_n$ only permutes coordinates, then
\begin{equation}\label{nd:multiplicative_wreath:eq:pure-coordinate}
 \boxed{[m_{(1^s)}]=
 \binom{n+|H|^{s-1}-1}{|H|^{s-1}-1}.}
\end{equation}
\end{theorem}

\begin{proof}
The fixed equation $a_rp^r(x)=x$ is equivalent to
$a_r=L_{p^r}(x)$.  Every nonempty fiber of the Lang map is a coset of
$\Fix_{N_n}(p^r)$, proving \eqref{nd:multiplicative_wreath:eq:lang} and the conditional spike law.
Averaging its $s$th power gives the middle expression in
\eqref{nd:multiplicative_wreath:eq:twisted-moment}; Burnside's lemma gives the orbit count.  Under
coordinate permutations, $N_n^{s-1}$ is an array of $n$ elements from
$H^{s-1}$ modulo column permutation, hence a multiset of size $n$ from
$|H|^{s-1}$ types.  This proves \eqref{nd:multiplicative_wreath:eq:pure-coordinate}.
\end{proof}

For coordinate permutations,
\[
 \Pr(F_1>0)=\frac{(1/|H|)_n}{n!}
 \sim\frac{n^{1/|H|-1}}{\Gamma(1/|H|)},
\]
again giving vanishing fixed points in probability with divergent moments.

\subsection{Varying base groups and the additive contrast}

Let $H_N$ act on $\Omega_N$ and $V_N$, and put
$G_{N,k}=H_N\wr S_k$.  Define
\[
 R_{N,s}=\#(H_N\backslash\Omega_N^s),
 \qquad a_{N,s}=\dim(V_N^{\otimes s})^{H_N}.
\]
Theorems~\ref{nd:multiplicative_wreath:thm:product-moment} and \ref{nd:multiplicative_wreath:thm:tensor} give the exact
two-parameter laws
\begin{align}
 [m_{(1^s)}]\MGF_{G_{N,k},\Omega_N^k}
 &=\binom{k+R_{N,s}-1}{R_{N,s}-1},\label{nd:multiplicative_wreath:eq:vary-product}\\
 [m_{(1^s)}]\MGF_{G_{N,k},V_N^{\otimes k}}
 &=\begin{cases}
 \binom{k+a_{N,s}-1}{a_{N,s}-1},&a_{N,s}>0,\\
 0,&a_{N,s}=0.
 \end{cases}\label{nd:multiplicative_wreath:eq:vary-tensor}
\end{align}
These binomial coefficients govern every joint-growth path.

The direct-sum block module behaves differently.  If
$A_N(\mathbf t)=\MGF_{H_N,V_N}(\mathbf t)$, the labelled-cycle exponential
formula gives
\begin{equation}\label{nd:multiplicative_wreath:eq:block-series}
 \sum_{k\ge0}u^k\MGF_{H_N\wr S_k,V_N^{\oplus k}}
 =\exp\!\left[\sum_{\ell\ge1}\frac{u^\ell}{\ell}
 A_N(t_1^\ell,t_2^\ell,\ldots)\right].
\end{equation}
Separating its constant term $(1-u)^{-1}$ shows coefficientwise that
\begin{equation}\label{nd:multiplicative_wreath:eq:block-limit}
 \lim_{k\to\infty}\MGF_{H_N\wr S_k,V_N^{\oplus k}}
 =\exp\!\left[\sum_{\ell\ge1}\frac{
 A_N(t_1^\ell,t_2^\ell,\ldots)-1}{\ell}\right].
\end{equation}
Thus the additive block representation has a compound-Poisson limit, while
\eqref{nd:multiplicative_wreath:eq:vary-product} and \eqref{nd:multiplicative_wreath:eq:vary-tensor} usually diverge.

\begin{corollary}[multiplicative wreath obstruction theorem]
For product and tensor-induced wreath representations, the coefficient
$[m_{(1^s)}]$ is a multiset count
$\binom{n+r_s-1}{r_s-1}$.  If $r_s\ge2$, it grows as
$n^{r_s-1}/(r_s-1)!$, so the raw $\sigma$-MGF has no finite coefficientwise
limit.
\end{corollary}

\subsection{Exact matrix verification}

\subsubsection{Independent computational routes}

The companion marimo notebook and \texttt{verification.py} use the following
audit protocol.
\begin{enumerate}
\item Construct every group element with the semidirect-product law.  Store a
permutation matrix by its unique nonzero row in each column; store a signed
monomial matrix by that row and a sign.  This is the complete matrix, not a
character-only surrogate.
\item Form matrix powers directly and compare their fixed points or traces
with the structural cycle-product, class-power, diagonal-centralizer, or Lang
map formula.
\item Compute $\bigotimes_j\Sym^{\tau_j}W$ on its actual monomial basis.
Permutation modules use literal configuration orbits.  Signed modules
propagate signs around each basis orbit: an orbit contributes one invariant
unless a stabilizer sends a basis vector to its negative.
\item Independently compute the same coefficient by the full Reynolds--Molien
average.  Moment cases receive a third comparison with the closed binomial or
centralizer formula.
\end{enumerate}
All checks use integers or exact rational numbers.

\subsubsection{Wreath types and product actions}

\begin{center}
\begin{tabular}{lrrrr}
\toprule
group & order & wreath types & power checks & weight checks\\
\midrule
$C_2\wr S_3$ & 48 & 10 & 288 & 10\\
$S_3\wr S_2$ & 72 & 9 & 432 & 9\\
\bottomrule
\end{tabular}
\end{center}

For $S_3\curvearrowright\{0,1,2\}$, direct product-action matrices give:
\begin{center}
\begin{tabular}{rrrrrr}
\toprule
$n$ & matrix dimension & $s$ & $R_s$ & direct/Burnside & formula\\
\midrule
2 & 9  & 2 & 2 & 3  & 3\\
2 & 9  & 3 & 5 & 15 & 15\\
3 & 27 & 2 & 2 & 4  & 4\\
3 & 27 & 3 & 5 & 35 & 35\\
\bottomrule
\end{tabular}
\end{center}
The program additionally checks all first six powers of all $72+1296$
product-action matrices, for $8{,}208$ fixed-power comparisons.

\subsubsection{Tensor-induced and fixed-tail signed groups}

Take $H=C_2$ and $V=\mathbf1\oplus\mathrm{sgn}$.  Then
$a_s=2^{s-1}$.  The tensor-induced results are:
\begin{center}
\begin{tabular}{lrrrrrr}
\toprule
group & dim. & $a_{(1)}$ & $a_{(2)}$ & $a_{(1,1)}$ & $a_{(2,1)}$ & $a_{(1^4)}$\\
\midrule
$C_2\wr S_2$ & 4 & 1 & 3 & 3 & 7  & 36\\
$C_2\wr S_3$ & 8 & 1 & 4 & 4 & 13 & 120\\
\bottomrule
\end{tabular}
\end{center}
Every entry is simultaneously a direct signed-fixed-space count and a Molien
coefficient; the moment entries also equal \eqref{nd:multiplicative_wreath:eq:tensor-moment}.  Traces
of six powers give another $336$ checks.

For the fixed-tail bipartition $((n-1),(1))$ of $C_2\wr S_n$, the natural
signed module has character polynomial $X_{1,+}-X_{1,-}$.  For $n=2,3,4$,
both direct and class-formula routes give
\begin{center}
\begin{tabular}{c|rrrrrr}
\toprule
$n$ & $(1)$ & $(2)$ & $(1,1)$ & $(3)$ & $(2,1)$ & $(1^4)$\\
\midrule
2 & 0&1&1&0&0&4\\
3 & 0&1&1&0&0&4\\
4 & 0&1&1&0&0&4\\
\bottomrule
\end{tabular}
\end{center}
The matrix trace also agrees with the power-transformed character polynomial
in $1{,}760$ element-power comparisons.  Its limiting character is the
Skellam variable $Z_+-Z_-$ with independent
$Z_\pm\sim\operatorname{Poisson}(1/2)$, concretely illustrating
\eqref{nd:multiplicative_wreath:eq:fixed-tail-limit}.

\subsubsection{Diagonal, twisted, Fock, and varying-base checks}

For the diagonal coset actions, matrix averages and
\eqref{nd:multiplicative_wreath:eq:diagonal-moment} agree:
\begin{center}
\begin{tabular}{lrrrr}
\toprule
action & dimension & $\E X$ & $\E X^2$ & $\E X^3$\\
\midrule
$S_3^2\curvearrowright S_3^2/\Delta S_3$ & 6  & 1 & 3  & 11\\
$S_3^3\curvearrowright S_3^3/\Delta S_3$ & 36 & 1 & 11 & 251\\
\bottomrule
\end{tabular}
\end{center}
All first six powers of all $252$ matrices give $1{,}512$ further checks.
For $S_3^2\rtimes S_2\curvearrowright S_3^2$, the direct/formula moments are
$1,21,666$ for $s=1,2,3$; all $432$ fixed-power checks satisfy the Lang-map
criterion, and the direct pair-orbit count is $21$.

For $C_2$ on $\mathbf1\oplus\mathrm{sgn}$, explicit bosonic monomials give
invariant dimensions $1,1,2,2,3,3,4$ in energies $0$ through $6$.
Explicit exterior bases give $1,1,0,0,0,0,0$.  These values and all
unaveraged power characters agree with \eqref{nd:multiplicative_wreath:eq:fock-plus} and
\eqref{nd:multiplicative_wreath:eq:fock-minus} in $94$ checks.  Regular bases $C_2,C_3$, with
$k=2,3$ and $s=2,3$, give eight further exact specializations of
\eqref{nd:multiplicative_wreath:eq:vary-product}--\eqref{nd:multiplicative_wreath:eq:vary-tensor}.

\resultbox{\textbf{Verification result.}  The executable suite reports
\textbf{PASS: 13,183 exact checks}.  This includes group construction,
semidirect powers, class weights, fixed points, traces, orbit counts, signed
fixed spaces, Molien coefficients, moment theorems, and graded Fock
characters.  The asymptotic claims are deductions from the proved exact
formulas; they are not extrapolations from the small cases.}

\subsection{Remarks and asymptotic classification}

\begin{center}
\small
\begin{tabularx}{\linewidth}{>{\raggedright\arraybackslash}p{0.25\linewidth}
>{\raggedright\arraybackslash}p{0.29\linewidth}X}
\toprule
family & controlling statistic & behavior\\
\midrule
$V^{\oplus n}$ & cyclewise determinants & compound-Poisson limit\\
$\Omega^n$ & $R_s=\#(H\backslash\Omega^s)$ & polynomial divergence if $R_s\ge2$\\
$V^{\otimes n}$ & $a_s=\dim(V^{\otimes s})^H$ & polynomial divergence if $a_s\ge2$\\
bosonic/fermionic Fock & $\det(I\mp q^rg^r)^{\mp1}$ & fixed-energy stability\\
fixed-tail multipartition & marked character polynomial & marked-Poisson chaos\\
growing multipartition & wreath Kronecker multiplicities & no universal raw limit\\
$H^n/\Delta H$ & common conjugacy class & exponential moment divergence\\
$H^n\rtimes P_n\curvearrowright H^n$ & $|\Fix_{H^n}(p)|$ & Bernoulli-spike rare events\\
varying $H_N\wr S_k$ & $R_{N,s}$ or $a_{N,s}$ & exact two-parameter criterion\\
\bottomrule
\end{tabularx}
\end{center}

The computational examples are deliberately modest: their role is to test
the representation, conventions, power operation, and coefficient extraction
without relying on asymptotic numerics.  The formulas themselves are valid at
every finite size covered by their hypotheses.  The small groups already
display the structural distinction: cyclewise sums exponentiate into stable
laws; cyclewise products create rare, large contributions that force moments
to grow.

\subsection{Reproducibility and scope}
\thispagestyle{fancy}

The standalone source, compiled PDF, marimo notebook, exact backend, and
focused tests live together in
\texttt{new\_dists/multiplicative\_wreath\_sigma\_mgfs}.  From the repository
root, the complete audit is reproduced with
\begin{center}
\texttt{.venv/bin/python -m
new\_dists.multiplicative\_wreath\_sigma\_mgfs.verification},
\end{center}
and the interactive laboratory is opened with
\begin{center}
\texttt{.venv/bin/marimo run
marimo\_notebooks/multiplicative\_wreath.py}.
\end{center}
The PDF is compiled directly from its colocated \LaTeX{} source.  No combined
anthology source or PDF is modified by this package.

The exhaustive computations certify the listed finite instances and test the
implementation of each structural reduction.  Universality comes from the
proofs: Reynolds averaging, wreath class powers, orbit--multiset bijections,
and the Lang-map fiber argument.  Statements about growing multipartitions
remain criteria rather than unconditional convergence theorems, because the
relevant wreath Kronecker multiplicities depend on the chosen family.

\begin{thebibliography}{9}
\bibitem{nd:multiplicative_wreath:wreath-symmetric}
I.~G.~Macdonald-style wreath characteristic background is developed in
\emph{Wreath Product Symmetric Functions},
\href{https://arxiv.org/abs/0809.2439}{arXiv:0809.2439}.

\bibitem{nd:multiplicative_wreath:fock}
I.~B.~Frenkel, N.~Jing, and W.~Wang,
\emph{Vertex representations via finite groups and the McKay correspondence},
\href{https://arxiv.org/abs/math/9907166}{arXiv:math/9907166}.

\bibitem{nd:multiplicative_wreath:fig}
K.~Casto,
\emph{$\mathrm{FI}_G$-modules, orbit configuration spaces, and complex
reflection groups},
\href{https://arxiv.org/abs/1608.06317}{arXiv:1608.06317}.
\end{thebibliography}
